\chapter{Системийн шинжилгээ ба зохиомж}
\section{Системийн ерөнхий архитектур}
Энэхүү судалгааны ажлын систем нь сошиал платформуудаас (Facebook, News, Blog) цугларсан өгөгдлийг хүлээн авч, BERT загвар болон NLP-д суурилсан аргуудаар боловсруулан харагдах байдалд оруулах архитектуртай. Үндсэн архитектур дараах хэсгүүдээс бүрдэнэ:
\begin{enumerate}
    \item \textbf{Data Collection (Өгөгдөл цуглуулах):} Сошиал эх сурвалжаас API эсвэл Scraper ашиглан Data Source үүсгэх.
    \item \textbf{Preprocessing (Урьдчилсан боловсруулалт):} Текст цэвэрлэх, зохисгүй үгс, emoji зэргийг цэвэрлэх алхам.
    \item \textbf{Embedding \& Topic Detection (Векторжуулалт ба Сэдэв илрүүлэх):} Mongolian BERT ашиглан текстийг хувиргах, дараа нь BERTopic эсвэл Cosine Similarity-д суурилсан сэдэв илрүүлэх.
    \item \textbf{NER болон Network Analysis (Нэрлэсэн объект таних болон сүлжээний шинжилгээ):} Сэдвийн дагуух хүмүүс, байршил болон байгууллагуудыг олох, тэдгээрийн уялдааг гаргах.
    \item \textbf{Sentiment Analysis:} Тухайн сэдэвт хэрхэн хандаж буйг эерэг, сөрөг, саармаг ангиллаар дүгнэх.
    \item \textbf{Dashboard (Хяналтын самбар):} Эцсийн үр дүнг хэрэглэгчид график, диаграмм (Network Graph, Bar charts) байдлаар харуулах.
\end{enumerate}

\section{Мэдээллийн урсгалын загвар (Data Flow)}
Системийн мэдээллийн урсгал нь "Data Source $\rightarrow$ Өгөгдөл цуглуулах модуль $\rightarrow$ Өгөгдлийн сан (Цуглуулсан өгөгдөл) $\rightarrow$ Өгөгдөл боловсруулалт ба шинжилгээ $\rightarrow$ Боловсруулсан өгөгдлийн сан $\rightarrow$ Хяналтын самбар" гэсэн дарааллаар ажиллана. 

\section{Тооцооллын болон системийн шаардлага}
\textbf{Моделийн хэмжээ ба Инференс хугацаа:} BERT base загварыг ашиглах үед 110M параметр буюу ойролцоогоор 420 MB хэмжээтэй байх бөгөөд T4 GPU ашиглан инференсийг ойролцоогоор 20-60 ms хугацаанд, секундэд 200-500 өгүүлбэр боловсруулах түвшинд тооцоолов. Харин CPU дээр 300-800 ms зарцуулдаг тул бодит production үед GPU болон боловсруулалтын нөөцийн зөв хуваарилалт шаардлагатай.

\section{NER Training ба Загварын сонголт}
NER сургахад \texttt{Davlan/bert-base-multilingual-cased-ner-hrl} HuggingFace загварыг ашиглах бөгөөд энэ нь өмнөх хуучирсан монгол хэлний нээлттэй загваруудын дутагдлыг (dependency conflict болон updated weights байхгүй) арилгаснаар Colab дээр шууд ажиллах боломжийг олгож байна.
